\section{Methodology}
\label{sec:method}
FISH currently targets ROS~2 Humble and is a tool chain with three parts: \emph{instrumentation}, tracepoints in the ROS~2 client libraries plus the CUDA profiler; \emph{acquisition}, which runs an unmodified application under that instrumentation and delimits the steady-state part of the trace; and \emph{extraction}, which turns the trace into the model of Section~\ref{sec:model}. The trace must answer four questions on one clock: which objects exist and when---each \texttt{rcl} handle from creation to destruction, so that a callback is attributed correctly even if a handle is reused; when each callback ran, whatever triggered it; which callback produced which message, including deliveries that never leave the process; and what GPU work each callback submitted, down to the kernels, copies and synchronisations.

\textbf{Instrumentation:} FISH builds on \texttt{ros2\_tracing}~\cite{bedard-2022-ros2tracing}, whose initialisation events record each object as it is created and whose execution events record what then happens to it, linked by the object's handle. To the 14 stock events of Humble, FISH adds 28. Finalize events close the object lifetimes that the stock init events open, and a timer event binds a timer to its node and period. Start and end events for waitables and clients, with client round-trip and action-server events, complete the five kinds of callback the executor dispatches. Link events at the publish, receive and client-request sites tie each message to the callback that issued it, and the intra-process subscription waitable covers deliveries that never reach \texttt{rcl}. Registration events give rclpy nodes the same callback chain as rclcpp, and introspection events record each executor's type and thread count and each callback's group. GPU work is recorded by CUPTI through \texttt{nsys}, aligned to the LTTng clock at ingest and bound through correlation ids to the API call and the callback that issued it.

\textbf{Acquisition:} The application is not modified: the instrumented libraries are installed on the host or in a container, and the FISH entry point wraps the \texttt{ros2 launch} command. Because CUPTI sees only processes started under the profiler, the entry point reads the launch description and starts the containers that will host GPU components under \texttt{nsys} from the outset, since composable nodes are loaded into a container only once. Standalone GPU processes are caught by the FISH daemon, which polls \texttt{/proc} and restarts any process holding a CUDA context under the profiler. Once the node list has been unchanged over five polls and every component is loaded, the daemon starts the input: a rosbag replay, or a schedule of timed commands for a live system. Phase boundaries come from the trace, so comparisons use the steady state: initialisation ends when every periodic timer has fired once, warm-up when every recurring callback has completed its first execution.

\textbf{Extraction:} Initialisation events create the vertices of Section~\ref{sec:model} and the edges the wiring declares; execution events add the edges only running reveals, with their measurements. The graph is checked against what the traced binaries create---a node constructor must yield one node, seven entities and seven callbacks---and is complete only if every such expectation appears in it.
